\subsection{我们可以为之辩护的内容}

在当前阶段，最能站得住脚的主张是定性的，而不是精确的：
\begin{enumerate}[nosep]
  \item 当当前状态的强化更强时，障碍往往会上升；
  \item 当目标状态更含糊时，障碍往往会上升；
  \item 当相关准备已经提前完成时，障碍往往会下降；
  \item 障碍既取决于当前条件，也取决于转变历史。
\end{enumerate}

\begin{definition}[准备与含糊性]
令 $U(s_1)$ 表示\textbf{目标含糊性}，即进入 $s_1$ 时仍未解决的规格负担。令 $P$ 表示\textbf{准备}，即在转变开始之前，与进入相关的工作已经提前完成到何种程度。
\end{definition}

\begin{remark}
\textbf{单调障碍支架。}
出于建模目的，可以把转变难度看作依赖于函数
\[
  \DV(s_0 \to s_1) \approx g\!\bigl(\widehat{\FI}(s_0), U(s_1), P, H\bigr),
\]
其中 $H$ 总结与转变相关的历史，而 $g$ 被假定为随着当前状态阻力与含糊性上升而上升，随着准备增加而下降。
\end{remark}

\begin{discussion}
这个假设比本章早先版本中的障碍分解公式更弱，也更容易辩护。它保留了后续章节所需要的业务逻辑：如果当前状态被强力强化、目标状态又很模糊、并且准备不足，那么转变就会更难。但它不再假装手稿已经识别出一条带有已知系数的特定加性定律。
\end{discussion}

\begin{remark}
这里的符号“$\approx$”应当被保守理解。它并不意味着统计拟合或渐近逼近。它只意味着：这个障碍构造正被表示为对这些已命名输入的某种建模依赖。
\end{remark}